\exercisesection{}

\begin{enumerate}
\def\labelenumi{\arabic{enumi})}
\item
  设 E 和 F 是巴拿赫空间，u 是一个连续线性映射 \(E \rightarrow F\) 。我们有 \(((u)) > 0\)，当且仅当 u 的像在 F 中是闭的。
\item
  设 \(p \in [1, +\infty]\) ，且 \(\mathrm{E} = \ell_{\mathbf{C}}^{p}(\mathbf{N})\) 。设 u 是 E 的一个对角自同态。对每个 \(\lambda \in C\) ，用 u 的谱计算 \(((u - \lambda 1_{\mathrm{E}}))\) 。
\item
  设 \(p \in [1, +\infty]\) 且 \(\mathrm{E} = \ell_{\mathrm{K}}^{p}(\mathbf{N})\)，并设 \(u\) 是移位自同态
\end{enumerate}

\[
(x _ {0}, x _ {1}, \dots ,) \mapsto (0, x _ {0}, x _ {1}, \dots)
\]

（E 上的）。证明 u 是指标为 1 的弗雷德霍姆自同态，且 \(((u)) = 1\) 。

\begin{enumerate}
\def\labelenumi{\arabic{enumi})}
\setcounter{enumi}{3}
\item
  设 E 是巴拿赫空间。映射 \(u \mapsto ((u))\)（从 \(\mathcal{L}(\mathrm{E})\) 到 R）是连续的，当且仅当 E 的维数 \(\leqslant 1\) 。
\item
  设 E 是希尔伯特空间，\(u \in \mathcal{L}(\mathrm{E})\) 。我们有 \(((u)) = \inf(\mathrm{Sp}(|u|) - \{0\})\)（\(|u|\) 的定义参见 I, p. 139, 定义 3）。
\item
  设 \(E = \ell^{2}(\mathbf{N})\)，\((e_{n})_{n \in \mathbf{N}}\) 是 E 的规范正交基。设 u 是 E 的对角自同态，满足 \(u(e_{0}) = e_{0}\) 且 \(u(e_{n}) = 2e_{n}\)（若 \(n \geqslant 1\)）。证明 \(((u)) = 1\)，且 \(u + h\) 对每个满足 \(\|h\| < 2\) 的 h 是弗雷德霍姆自同态。
\item
  设 E 和 F 是巴拿赫空间，u 是 \(\mathcal{Q}\mathcal{M}(\mathrm{E};\mathrm{F})\)（相应地，\(\mathcal{Q}\mathcal{E}(\mathrm{E};\mathrm{F})\)）的一个元素。设 h 是 \(\mathcal{L}(\mathrm{E};\mathrm{F})\) 的一个元素。假设存在正实数 a 和 b，使得 \(a + b((u))^{-1} < 1\)，从而
\end{enumerate}

\[
\| h (x) \| \leqslant a \| u (x) \| + b \| x \|\tag{2}
\]

对所有 \(x \in E\) 成立。

证明 \(u + h \in \mathcal{Q}\mathcal{M}(\mathrm{E};\mathrm{F})\)（相应地，\(u + h \in \mathcal{Q}\mathcal{E}(\mathrm{E};\mathrm{F})\)）且

\[
\begin{array}{c} \dim \operatorname{Ker} (u + h) \leqslant \dim \operatorname{Ker} (u), \\ \dim \operatorname{Coker} (u + h) \leqslant \dim \operatorname{Coker} (u), \\ \operatorname{ind} (u + h) = \operatorname{ind} (u). \end{array}
\]

（化归到 \(a\) 和 \(b\) 都 \(>0\) 的情形，并设 \(c > 0\) 满足 \(a + b/c < 1\) ；考虑 E 上由下式定义的范数 \(p\)

\[
p (x) = a \| u (x) \| + b \| x \| \quad \text {   for   } x \in \mathrm{E},
\]

并设 \(E_{1}\) 是赋有范数 p 的巴拿赫空间 E；我们有 \(\mathcal{Q}\mathcal{M}(E_{1};F)=\mathcal{Q}\mathcal{M}(E;F)\)，\(\mathcal{Q}\mathcal{E}(E_{1};F)=\mathcal{Q}\mathcal{E}(E;F)\) 且 \(\mathcal{L}(E_{1};F)=\mathcal{L}(E;F)\) 。然后应用 III, p. 62 的引理 2。）

\begin{enumerate}
\def\labelenumi{\arabic{enumi})}
\setcounter{enumi}{7}
\item
  设 E 是希尔伯特空间。不用博苏克–乌拉姆定理，直接证明 III, p. 64 的定理 4。
\item
  设 \(r \geqslant 1\) 是整数。若 \(f: \mathrm{M}_1 \to \mathrm{M}_2\) 是 \(\mathrm{C}^r\) 类流形之间的一个 \(\mathrm{C}^r\) 类态射，则称 \(x \in \mathrm{M}_1\) 是 \(f\) 的临界点，如果 \(f\) 在 \(x\) 处不是淹没（VAR, R1, 5.9.1）。设 X 是 \(f\) 的临界点集。集合 \(f(\mathrm{X})\)（在 \(\mathrm{M}_2\) 中）称为 \(f\) 的临界值集，而 \(f(\mathrm{M}_1) = f(\mathrm{X})\) 称为 \(f\) 的正则值集。
\end{enumerate}

若 \(M_{1}\) 和 \(M_{2}\) 是 \(C^{r}\) 类实流形（VAR, R1, 5.1.5, p. 35），流形的一个态射 \(f: M_{1} \to M_{2}\)（\(C^{r}\) 类）称为弗雷德霍姆的，如果对每个 \(x \in M_{1}\) ，切线性映射 \(\mathrm{T}_{x}(f)\) 是从巴拿赫空间 \(\mathrm{T}_{x}(\mathrm{M}_{1})\) 到巴拿赫空间 \(\mathrm{T}_{f(x)}(\mathrm{M}_{2})\) 的弗雷德霍姆自同态。a) 设 \(f: M_{1} \to M_{2}\) 是一个弗雷德霍姆态射。若 \(M_{1}\) 是连通的，则指标 \(\operatorname{ind}(\mathrm{T}_{x}(f))\) 不依赖于 \(x \in M_{1}\) 。它称为弗雷德霍姆态射 f 的指标，记作 \(\operatorname{ind}(f)\) 。

\begin{enumerate}
\def\labelenumi{\alph{enumi})}
\setcounter{enumi}{1}
\tightlist
\item
  设 U 是巴拿赫空间 E 的一个开子集，满足 \(0 \in U\)，并设 F 是巴拿赫空间。设 \(f: U \to F\) 是一个 \(C^{r}\) 类弗雷德霍姆态射，并设 \(u = \mathrm{T}_{0}(f)\) 。存在一个巴拿赫空间 \(E_{1}\)、U 的一个卡 \((\mathrm{U}_{1}, \varphi, \mathrm{E}_{1} \times \mathrm{Ker}(u))\)（满足 \(0 \in U_{1}\)）、F 的一个卡 \((\mathrm{U}_{2}, \psi, \mathrm{E}_{1} \times \mathrm{Coker}(u))\)（满足 \(f(0) \in U_{2}\)），以及一个弗雷德霍姆态射 \(g: \varphi(\mathrm{U}_{1}) \to \mathrm{E}_{1} \times \mathrm{Coker}(u)\)（\(C^{r}\) 类），使得
\end{enumerate}

\[
(\psi \circ f \circ \varphi^ {- 1}) (a, b) = (a, g (a, b))
\]

对所有 \((a,b)\in \varphi (\mathrm{U}_1)\subset \mathrm{E}_1\times \mathrm{Ker}(u).\) 成立。

\begin{enumerate}
\def\labelenumi{\alph{enumi})}
\setcounter{enumi}{2}
\item
  设 \(f \colon \mathrm{M}_1 \to \mathrm{M}_2\) 是 \(\mathbf{C}^r\) 类流形之间的一个 \(\mathbf{C}^r\) 类弗雷德霍姆态射。对每个 \(y \in \mathbf{M}_2\)，存在一个开邻域 \(\mathrm{V}\)（\(y\) 的），使得 \(f\) 通过过渡到子空间定义一个逆紧映射 \(f^{-1}(\mathrm{V}) \to \mathrm{V}\) 。
\item
  设 U 是巴拿赫空间 E 的一个包含 0 的开子集，并设 F 是巴拿赫空间。设 \(f \colon \mathrm{U} \to \mathrm{F}\) 是一个 \(\mathrm{C}^{\infty}\) 类弗雷德霍姆态射。对每个 \(x \in \mathrm{U}\) 和 F 中 \(f(x)\) 的每个邻域 V，存在 \(f\) 在 V 中的一个正则值。（利用第二萨德定理，VAR, R2, 10.1.3, p. 37。）e) 设 \(f \colon \mathrm{M}_1 \to \mathrm{M}_2\) 是 \(\mathrm{C}^{\infty}\) 类流形的一个 \(\mathrm{C}^{\infty}\) 类弗雷德霍姆态射。若 \(\mathrm{M}_1\) 和 \(\mathrm{M}_2\) 具有可数基，则 \(f\) 的临界值集是贫集（“斯梅尔定理”）。
\end{enumerate}

\(f)\) 设 \(f\colon \mathrm{M}_1\to \mathrm{M}_2\) 是具有可数基的 \(\mathrm{C}^{\infty}\) 类流形的一个 \(\mathrm{C}^{\infty}\) 类弗雷德霍姆态射。若 \(\operatorname {ind}(f) < 0\)，则 \(f\) 的像的内部是空的。

\(g)\) 设 \(f\colon \mathrm{M}_1\to \mathrm{M}_2\) 是具有可数基的 \(\mathrm{C}^{\infty}\) 类流形的一个 \(\mathrm{C}^{\infty}\) 类弗雷德霍姆态射。对每个位于一个贫集之外的 \(y\in \mathrm{M}_2\)，集合 \(f^{-1}(y)\) 是维数为 \(\operatorname {ind}(f)\) 的纯子流形。
% semantic-fragments: legacy-input-seam
\relax{}
